% mage-auto-heading: preserve-engineering-freedom
% mage-auto-heading: catalogue-the-controls
% mage-auto-heading: model-interactions-between-controls
% mage-auto-heading: trace-obligations-to-mechanisms
% mage-auto-heading: add-reconcile-and-retire-controls
% Authoritative LaTeX manuscript; Markdown is a generated compatibility projection.

% <!-- part-title: Alignment -->

% <!-- chapter-title: When Controls Become a System -->

\chapter{When Controls Become a System}

% <!-- gloss-only: Governance graph | A structured model whose nodes are the fleet's own governance mechanisms (each tagged by the event it fires on and the shared resources it reads, writes, or locks) and whose edges are the conflicts between them over a shared resource, in a closed taxonomy of contradiction, contention, ordering, and soft-versus-hard; the dynamic dual of a static catalogue of controls, holding the interaction dimension a list omits. -->

% <!-- point: every-control-reasonable-alone-the-system-of-controls-still-wrong | Every control can be reasonable in isolation while the control machinery as a whole remains inconsistent, incomplete, or obsolete. | terms: governance-graph, constraint -->

As governance mechanisms accumulate, they become control machinery: the constraints, sensors, and gates through which the environment governs agentic work. The machinery can fail even when every individual control is reasonable. Two controls can impose incompatible demands; their execution order can be ambiguous; an obsolete control can survive after its purpose disappears; an important obligation can remain uncovered. None of these failures is visible by reading one mechanism\textquotesingle s code.

% <!-- point: successful-environments-accumulate-machinery-so-the-estate-is-a-portfolio-to-manage | Governance conversion and up-front design accumulate control machinery; that machinery becomes an engineering object to maintain, reconcile, and sometimes retire. | terms: governance-graph, engineering-capital -->

% <!-- point: alignment-turned-on-itself-a-model-of-the-controls-made-binding-on-the-controls | Alignment turned on itself: a model of the controls, enforced on the controls. | terms: governance-graph -->

These interactions require engineers to maintain the controls as a system. MAGE asks three questions: \textbf{What controls exist? How do they interact? Which obligations do they cover?} These questions produce three useful views: a catalogue, an interaction graph, and a coverage relation. Represent that machinery explicitly, so engineers and tools can reason about how its controls compose and what obligations they cover---and, once the machinery materially affects engineering work, hold it to obligations of its own.

\section{Catalogue the Controls}\label{heading:catalogue-the-controls}

% <!-- section-terms: governance-graph -->

% <!-- point: the-catalogue-gives-each-control-a-stable-identity | Catalogue each control: a stable identity, the boundary it acts at, the resources it touches, the obligation it serves, and its governance role. | terms: governance-graph, shared-resource -->

Start by giving each control a stable identity: the event or boundary at which it acts, the resources it reads or changes, the obligation or failure class that justifies it, and its governance role or roles. The catalogue supplies stable identities for the interaction and coverage views, so engineers and tools need not rediscover controls from scripts, hooks, CI files, and institutional memory.

% <!-- point: a-catalogue-alone-cannot-show-conflict-or-a-missing-control | A catalogue can describe each control individually, but it cannot reveal conflicts between them; conflicts appear in the relationships among controls. | terms: governance-graph, shared-resource -->

A catalogue can describe each control individually, but it cannot reveal conflicts between them. Those conflicts appear in the relationships among controls.

\section{Model Interactions Between Controls}\label{heading:model-interactions-between-controls}

% <!-- section-terms: governance-graph, shared-resource -->

% <!-- point: each-collision-lives-in-the-space-between-two-controls-drawn-nowhere | A collision lives between two controls, drawn nowhere until the graph draws it. | terms: shared-resource, governance-graph -->

Suppose one repository check requires a task's changes to be squashed into one commit, while another rejects a commit whose diff exceeds a prescribed scope. A task may satisfy the first rule only by violating the second. Inspecting either check alone will not reveal how their requirements interact. \autoref{collision-edge} draws the shared commit set on which both act. The engineer must examine that relation to decide whether the obligations conflict or the mechanisms need a different arrangement.

\begin{figure}[htbp]
\centering
\includegraphics[width=\linewidth]{book/latex/assets/collision-edge-fe5071601b.pdf}
\caption{\emph{The conflict is the edge.} Two controls may be individually valid yet impose incompatible demands on the same resource. The interaction belongs to the relation among the controls and shared resource, not necessarily to either control\textquotesingle s implementation.}
\label{collision-edge}
\end{figure}

% <!-- point: the-interaction-graph-joins-controls-by-shared-resource | The interaction graph uses controls as nodes and shared resources as the principal join key; controls need not call one another to collide. | terms: governance-graph, shared-resource -->

The interaction graph uses controls as nodes and shared resources as the principal join key. A resource may be ordinary --- a lock, a file, a commit-set, a deployment slot --- or operational, such as a lifecycle event or a context budget. Controls need not call one another to interact; they only need to make potentially incompatible demands on the same modeled resource.

% <!-- point: shared-resource-is-a-useful-conservative-relation-not-a-proof-of-completeness | Shared resource is a useful conservative relation, not a proof that the graph captured every interaction. | terms: shared-resource, semantic-gap -->

A shared resource is a conservative signal of interaction, not proof that every conflict has been captured. Within the represented surface, it lets engineers and tools ask useful questions before the controls interact in production work: which controls fire on this event, which touch this resource, which pairs have no declared order, and which known combinations impose incompatible requirements.

% <!-- point: decidable-edges-a-lint-settles-semantic-ones-a-validator-or-human-judges | Some candidate conflicts are mechanically decidable; others require judgment about what the mechanisms demand. | terms: governance-graph, semantic-gap -->

Some candidate conflicts can be evaluated mechanically, such as a cycle in declared ordering requirements. Others need judgment. A static check can show that two controls touch the same commit set; a validator or human reviewer must determine whether their requirements conflict. Producing evidence of an interaction is different from judging it unacceptable. The graph makes the pair available for evaluation before the controls interfere during work.

% <!-- point: diagnose-the-conflict-kind-before-weakening-an-obligation | Not every conflict between controls reflects a conflict between the obligations they enforce; incompatible obligations require revisiting requirements, while interfering mechanisms may be changed without weakening either obligation. | terms: governance-graph, shared-resource -->

Not every conflict between controls reflects a conflict between the obligations they enforce. Two obligations may genuinely be incompatible, requiring engineers to revisit the requirements or specification. But compatible obligations can also be enforced by mechanisms that interfere with one another. In that case, changing the mechanisms or how they interact may resolve the conflict without weakening either obligation.

% <!-- point: governance-interaction-graph-candidate-conflict-condition | A read/write relation identifies potential interactions at a shared event; classifying the conflict and resolving it are separate judgments. | terms: governance-graph, shared-resource, semantic-gap -->

A simple interaction model records each mechanism's execution event, or \emph{slot}, and the sets of resources it reads (\(R\)) and writes (\(W\)). Two mechanisms have a \textbf{potential conflict} when they execute at the same event and at least one writes a resource the other accesses. Two readers do not conflict under this relation. In a graph \(G = (V, E)\), an undirected edge joins distinct mechanisms \(u\) and \(v\) when

\[
\operatorname{slot}(u) = \operatorname{slot}(v) \land ((W_u \cap R_v) \cup (R_u \cap W_v) \cup (W_u \cap W_v)) \neq \varnothing
\]

An edge identifies an interaction requiring attention; it does not establish that the underlying obligations are incompatible. The model must record the evaluation and any chosen resolution separately.

% <!-- point: three-turn-end-hooks-one-shared-slot | DocAble's worked control-system case connects a missing enforcement consequence to later modeling of interactions among three turn-end hooks. | terms: governance-graph, shared-resource -->

% <!-- box-family: inset -->

% <!-- inset-domain: Case study -->

% mage-box-prefix:

\begin{bookbox}{inset}{Worked inset --- DocAble: Governing Development Automation}
\label{inset-docable-control-system}
Controls become software that must itself be understood and maintained. DocAble's development environment assigns work to agents in separate Git worktrees, merges their changes, and runs checks around the agents' work. That automation creates obligations of its own: preserve contributions during integration, recognize unfinished work, and avoid having one automation prevent another from doing its job. Two episodes show why evaluating each script in isolation was insufficient.

In the first, merge automation processed a backlog of worktrees. A convenience option resolved conflicts by automatically taking one side, which silently dropped content when edits overlapped. An audit already detected conflicts affecting files that described or controlled the engineering environment and correctly judged automatic resolution unsafe. It only logged a warning, so the merge continued. The repair removed the unsafe fallback and sent ambiguous merges to a human. \autoref{judgment-without-consequence} shows the missing connection between the correct judgment and the permitted operation.

\begin{figure}[htbp]
\centering
\includegraphics[width=\linewidth]{book/latex/assets/judgment-without-consequence-6c9128b4cf.pdf}
\caption{\emph{A correct warning did not prevent loss.} The merge audit detected the dangerous conflict, but no mechanism acted on its verdict. The repair removed automatic resolution for ambiguous merges and required human involvement.}
\label{judgment-without-consequence}
\end{figure}

A different problem arose among hooks that run when an agent finishes a turn. One discourages premature termination, another prompts reflection on failures, and a third checks for stranded work. All three write to the turn's closing message. Each can be useful alone, but their order affects whether the others can complete their work. The partial governance graph examined here records twenty-five mechanisms and three conflict edges, all among these hooks. It models selected interactions whose cost justified explicit representation, rather than claiming to cover every possible conflict. \autoref{turn-end-ordering} shows the order declared for these three hooks.

\begin{figure}[htbp]
\centering
\includegraphics[width=\linewidth]{book/latex/assets/turn-end-ordering-8b921a2d5a.pdf}
\caption{\emph{Three hooks share the closing message.} The model orders the premature-termination guard, reflection, and the stranded-work check. It computes the order and reports cycles; the generic runtime driver needed to execute that computed order has not yet been implemented.}
\label{turn-end-ordering}
\end{figure}

The graph classifies conflicts and records a resolution such as merging mechanisms, ordering them, separating their execution points, accepting the interaction, or leaving it open. For these three hooks, pairwise ordering constraints put the premature-termination guard first, reflection second, and the stranded-work check last. Topological sorting computes the sequence and reports a cycle as an error. A separate check reports discrepancies between the governance model and the configured mechanisms. The generic driver that would execute hooks in the computed order had not yet been built, so this ordering remained an explicit, checkable design decision rather than an automatically enforced sequence.

These episodes lead to different repairs. The merge failure required removing a dangerous capability; the hook interactions required an explicit account of shared resources and ordering. A catalogue of individual controls would not have supplied either conclusion by itself. Modeling the relevant behavior and interactions lets engineers decide what to remove, what to coordinate, and which accepted decisions still need an enforcing mechanism.
\end{bookbox}

\section{Trace Obligations to Mechanisms}\label{heading:trace-obligations-to-mechanisms}

% <!-- section-terms: governance-graph, governance-conversion -->

% <!-- point: a-catalogue-and-a-graph-do-not-say-what-should-be-governed-but-is-not | A catalogue says what exists and a graph says where controls collide; neither says what should be governed but is not. | terms: governance-graph, governance-conversion -->

Interactions are only half of governing the control machinery. A catalogue tells you what controls exist; an interaction graph tells you where controls may collide. Neither tells you what \emph{should} be governed but is not. For that, join the mechanisms to the obligations or failure classes that justify them.

% <!-- point: coverage-is-a-requirements-traceability-problem-gap-vs-orphan | Coverage is a small requirements-traceability problem: an obligation with no adequate mechanism is a coverage gap; a mechanism with no obligation is an orphan. | terms: governance-graph, governance-conversion -->

The result is a small requirements-traceability problem. An important obligation with no adequate mechanism is a coverage gap. A mechanism with no current obligation, failure class, or documented purpose is an orphan. The useful question is not "how many controls do we have?" It is whether consequential obligations have appropriate coverage and whether every surviving control can still explain why it exists. \autoref{census-trio} sets the three views side by side over the same control machinery.

\begin{figure}[htbp]
\centering
\includegraphics[width=\linewidth]{book/latex/assets/census-trio-ab9995bb33.pdf}
\caption{\emph{Three views of the control machinery.} The catalogue identifies controls; the interaction view relates controls through shared resources; the coverage view relates controls to the obligations that justify them. Stable identities connect the three views while letting each answer a different engineering question. A missing mechanism appears as a coverage gap; an untraced mechanism appears as a candidate orphan.}
\label{census-trio}
\end{figure}

% <!-- point: gap-is-under-governance-orphan-is-governance-debt-neither-automatically-bad | A coverage gap is an uncovered obligation; an orphan is a mechanism with no current purpose. Neither is automatically bad — an orphan is a finding to reconcile, not proof to delete. | terms: governance-graph -->

Two findings do the diagnostic work, and neither verdict is automatic.

\begin{itemize}
\tightlist
\item
  \textbf{Coverage gap.} A consequential obligation for which the intended evidence or enforcement is not adequately supplied.
\item
  \textbf{Orphan.} A mechanism that no longer traces to a current obligation, failure class, or documented purpose. An orphan is a finding to reconcile, not proof that the mechanism should be deleted.
\end{itemize}

\section{Add, Reconcile, and Retire Controls}\label{heading:add-reconcile-and-retire-controls}

% <!-- section-terms: governance-graph, teetering-tower, engineering-capital -->

% <!-- point: conversion-is-reactive-coverage-is-proactive-a-healthy-environment-needs-both | Governance conversion is reactive and coverage auditing is proactive; a healthy environment needs both. | terms: governance-graph, governance-conversion -->

Governance conversion and coverage auditing provide complementary pressures. Conversion adds durable structure when experience exposes a recurring gap; coverage auditing finds controls that are missing, redundant, or obsolete.

% <!-- point: the-same-audit-pushes-against-permanent-accumulation-capital-is-not-the-number-of-controls | The same audit pushes against permanent accumulation: engineering capital is not the number of controls. | terms: engineering-capital, teetering-tower -->

The same audit also pushes against permanent accumulation. A control whose purpose has disappeared should be reconciled or retired; two controls that now overlap may be consolidated; a mechanism whose maintenance cost exceeds the failure it prevents may no longer be worth carrying. This echoes \magevalue{sec:governance-conversion}: \textbf{engineering capital is not the number of controls.} A control that no longer produces value has depreciated into overhead; the task now is to reconcile or retire it.

% <!-- point: the-meta-model-stays-connected-where-controls-are-modeled-explicitly | Entries in a control model identify the implementing mechanism, and correspondence checks expose missing or unrecognized targets instead of silently treating them as covered. | terms: traceability, teetering-tower -->

The control model must remain connected to the controls it represents. An entry can identify the script or configured permission that implements a control, while a correspondence check reports missing or unrecognized targets. The check should expose a gap in the relation rather than silently treating an unmatched control as covered. This applies the broader rule from \magevalue{chapter:modeling}: a model of the engineering environment must stay connected to the machinery that engineers and agents use it to reason about.

% <!-- point: stop-the-recursion-unless-experience-earns-more-model-the-estate-that-exists | Stop the recursion unless experience earns more: a governance model can itself become bureaucracy. | terms: governance-graph, teetering-tower -->

Do not recursively govern every relation merely because it can be modeled. Extend the control model only when a recurring engineering question or observed failure justifies the added machinery.

The worked control-machinery model is grounded primarily in the originating case. Its broader value is a hypothesis for other engineering environments: the catalogue, interaction, and coverage views provide concrete questions to test when a growing set of controls becomes difficult to manage.

Governance conversion adds durable engineering structure; governing the control machinery keeps that structure coherent. Catalogue the controls, expose their interactions, trace them to the obligations they serve, and retire machinery that no longer earns its cost.

\subsection{Preserve Engineering Freedom}\label{heading:preserve-engineering-freedom}

% <!-- point: alignment-distinguishes-what-must-be-governed-from-what-may-remain-open | Alignment must preserve two complementary freedoms — freedom of implementation (which acceptable realizations remain available) and freedom to extend the environment (whether additional controls can coexist); the realization-space model answers the first, the governance graph the second. | terms: governance-graph, alignment-principle -->

Alignment must preserve two kinds of freedom. The first is freedom of \emph{implementation}: which acceptable realizations remain available after obligations are enforced. The second is freedom to \emph{extend the engineering environment}: whether additional controls can coexist with those already installed. The realization-space model in \magevalue{sec:constraints-sensors-validators-gates} addresses the first question; the governance graph addresses the second.

% <!-- point: governing-the-governors-is-second-order-not-a-license-for-unrestricted-agents | Modeling and Alignment extend to the machinery that performs engineering work; the result is not unrestricted agent autonomy — agents may explore without being authorized to commit, and engineers still decide which freedoms to delegate. | terms: governance-graph -->

This extends Modeling and Alignment to the machinery that performs engineering work. Product obligations motivate controls; controls interact; and those interactions must themselves be modeled and governed. The result is not unrestricted agent autonomy. Agents may be permitted to explore alternatives without being authorized to merge, deploy, or otherwise commit them. Engineers must still decide which freedoms can be delegated and which require stronger controls.



% <!-- point: modeling-and-alignment-are-both-on-the-table-but-that-is-not-yet-a-method | Chapter 4 shows how engineers use and revise models and controls together as an engineering problem develops. | -->

Engineers must decide when a question deserves a model, when an obligation deserves enforcement, how work should move through the environment, and what to do when the available representation or control is insufficient. Building on these accounts of Modeling and Alignment, \magevalue{chapter:one-problem-many-models} shows how engineers use and revise models and controls together as an engineering problem develops.
